\subsection{The target slope floor in the original comparison}
\label{geom:actual-target-floor}

Fix an actual translate and write \(R_0=RR\) for its row lattice,
\(R=RR^\vee\) for the target, and
\(s=s_{\rm Frame}=\langle\mu(N'_i)-h_i\rangle\) for the original
rank-weighted source centre. The formal data set
\(\mu_{\min,CC}=\mu_{\min}(R)\) and \(HH=x^{C_H}b\).
The exact dual-slope identity turns the required floor into
\[
 s-\mu_{\min}(R)=s+\mu_{\max}(R_0)\le x^{C_H}b.
\]
The duality statement uses the same field and the original dual norm.
\src{src/Formalism/LocalRingCardinality/DataArith.lean}{102}{the minimum of the dual lattice}
\src{src/Formalism/LocalRingCardinality/RankBoundMonotoneInDimension.lean}{993}{the exact dual-slope identity}
\src{src/Formalism/SharpRestriction/GlobalSchurBounds/QuadraticLogDerivCongr.lean}{714}{transfer into the rank-gap inputs}.

\paragraph{An auxiliary family of tests on all rows.}
The source obtains the floor by constructing a map out of the
\emph{primal} row space on the base level using the auxiliary height tests. In the original
\code{Separates} statement every nonzero row \(t\) in the specified
\(\mathrm{EEat}\) space has a nonzero pairing with some auxiliary form
\(\Psi_{\rm test}\) and some Schwartz datum \(\Phi\) in the grid model.
The quantifiers first choose the real vector and covector for the
given multiplicity vector. The witness then varies with \(t\).
The original separation theorem uses the real branching period,
definite-level theta surjectivity, Bessel ascent, and level-tower
convergence and gap conditions.
\src{src/Development/AbelianVarietyTheory/BlockTheory/DivisorBase/LevelPlusSubgroup.lean}{65}{auxiliary row-separation quantifiers}
\src{src/Development/AutomorphicArithmetic/ThetaCorrespondence/ThetaPairLifts/LocalTypeDictionary/RealFunctionalsSeparateOfHypothesis.lean}{15}{proof from the local theta and real inputs}.

The fitted grid replaces these tests by jointly injective coordinates
\(\psi_K\). The base-row map \(v_{\mu_0}\) has kernel
\(\ker P_q(\mu_0)\). Indeed the formula for each fitted coordinate
compares it with \(\psi_K\) of the theta image, with nonzero scalar
factors. The theta image itself is zero precisely when the base
quotient coordinate \(P_q(\mu_0)\) is zero. The formal proof uses the
surjectivity of \(P_q(\mu_0)\) and one nonzero row to certify the
common scalar before comparing kernels. The fitted maps also give
\[
 \Psi\,\iota_tP_q(t)=
          \left(\prod_{i,j}r_q(t)_{ij}\right)^{-1}v_t
\]
at the base tuple and every tuple occurring at a finite place, with
the factor equal to one at the base. It follows that \(\Psi\) is
injective. For if \(\Psi z=0\), lift \(z\) to
\(\iota_{\mu_0}P_q(\mu_0)y\). Then \(v_{\mu_0}y=0\), so
\(P_q(\mu_0)y=0\) and \(z=0\). These are auxiliary fitted tests on rows.
The construction provides the target floor without a rank assertion
about the actual word columns \(B_i\).
\src{src/Development/AutomorphicArithmetic/ThetaCorrespondence/ThetaPairLifts/ThirdPlaceSeries/FittedFunctionAtZeroClose.lean}{51}{use of row separation in the fitted function}
\src{src/Development/AutomorphicArithmetic/ThetaCorrespondence/ThetaPairLifts/ThirdPlaceSeries/KerEqOfFormula.lean}{20}{kernel of the fitted matrix}
\src{src/Formalism/LocalToGlobal/IndefiniteForms/IndefiniteTree/FunctionalLowerBoundOfConditionsNumeric.lean}{917}{injectivity of the induced row map}.

\paragraph{Local and infinite costs of the same map.}
Write \(R_{\mathrm b}\) for the base row lattice and \(M_{\mathrm{aux}}\)
for the auxiliary target of \(\Psi\). The finite-place comparison uses
the actual glued lattice
\(
R_{\mathrm b}=\operatorname{glue}_w
(\iota_{\mu(w)}\,\mathrm{abelian}_2(\mu(w)))
\).
The fitted local inclusion, with the correction \(d_F(w)\), gives
\[
 \Psi(R_{\mathrm b})_w\ \subseteq\
 \mathfrak p_w^{-a_w}(M_{\mathrm{aux}})_w,\qquad
 a_w=v_w\!\left(\prod_{i,j}r_q(\mu(w))_{ij}\right)+d_F(w).
\]
The support is the union of the nonbase-tuple places and the support
of \(d_F\). Local-to-global inclusion therefore gives
\(\Psi(R_{\mathrm b})\subseteq (\prod_w\mathfrak p_w^{-a_w})
M_{\mathrm{aux}}\). At infinity the source first rescales the row
frame by \(\kappa^{-1}\) via \(\code{scaleFrame}\), with its original
squared-norm convention. The resulting source maximum slope is
\(\mu_{\max}(R_{\mathrm b})+\operatorname{mean}\log\kappa\).
The archimedean operator height in this metric is bounded by
\(b_{\rm IP}\). The injective-map slope inequality and ideal scaling
then give the signed estimate
\begin{equation}\label{geom:auxiliary-row-height}
 \mu_{\max}(R_{\mathrm b})+\operatorname{mean}\log\kappa
 \le \mu_{\max}(M_{\mathrm{aux}})+b_{\rm IP}
                         +2\sum_w\lambda_wa_w,
 \qquad
 \lambda_w=\frac{\log N\mathfrak p_w}{[K:\mathbf Q]},
\end{equation}
where \(K\) is the number field of the lattices. The sign follows
from the negative power of \(\mathfrak p_w\) in the containing ideal.
There is no additional row-dimension multiplier. The original
exponent-grid bound charges the last sum to the grid and fixed-point
budgets. Together with the auxiliary lattice bound and
\(b_{\rm IP}\), it gives the input denoted \(\code{FUN}\).
The base-to-sharp passage adds \(14\lambda_0\) to that maximum-slope
bound, with \(\lambda_0\) the original logarithmic group-size term.
The metric and row-lattice ties identify the sharp maximum slope
with that of \(RR\).
\src{src/Formalism/LocalToGlobal/IndefiniteForms/IndefSide.lean}{281}{glued row lattice}
\src{src/Formalism/LocalToGlobal/IndefiniteForms/IndefiniteTree/FunctionalLowerBoundOfConditionsNumeric.lean}{973}{local-to-global inclusion and infinite-place estimate}
\src{src/Formalism/LocalToGlobal/ExponentGrids/ExponentGridsProofs/FrameHeightOfDualSlope.lean}{405}{slope inequality with fractional-ideal denominators}
\src{src/Development/AutomorphicArithmetic/ThetaCorrespondence/ThetaPairLifts/ThirdPlaceSeries/HefunPointB.lean}{32}{transport from the base to the actual sharp row lattice}.

\paragraph{The common normalization and the extreme tail.}
The normalization term \(\operatorname{mean}\log\kappa\) cancels
against the negative of that term in the word heights \(h_i\).
The other source terms are bounded through the upper source
maximum-slope comparison, the base column maximum slope, the gauge
combination \(G+F^0-E_i\), and the remaining height and rank budgets.
The rank-weighted average gives the required bound on \(s+\mu_{\max}(RR)\).
These bounds give a fixed polynomial exponent \(C_H\) carried into
the frame construction. They do not estimate the target floor on
the scale of the final error bracket by itself.
\src{src/Formalism/LocalRingCardinality/KernelPlaceMeanBound.lean}{25}{row scalar in the word height}
\src{src/Formalism/SharpRestriction/Divisibility/SlopeExcess/HeightBoundOfComponentBounds.lean}{13}{rank-weighted averaging of the source bounds}
\src{src/Formalism/SharpRestriction/Divisibility/SlopeExcess/HeightBoundOfComponentBounds.lean}{80}{assembly of component bounds}.

The frame estimate concerns the same word maps, the same row
lattice, and the same exponent. At the chosen infinite place the
full maps are compressed by isometries, and their retained rows
are the actual \(B_i\) in the row and column metric frames. The
source uses \(\sigma=x^{-C_H}\), a compression loss \(d_e\) and a
diagonal defect majorant \(e_B\) with
\[
 (3\sqrt{d_e}+4\sqrt{e_B})x^{C_H}\le1,\qquad
 \eta\le 5x^{-C_H}+3\sqrt{d_e},\qquad e\le e_B .
\]
Consequently, by monotonicity of the square root,
\[
       (\eta+4\sqrt e)x^{C_H}\le6,\qquad
       (\eta+4\sqrt e)HH\le6b .
\]
The tail-counting argument bounds the fraction below \(s-t\) by
\((1-\delta_{\rm row})^{-1}(\eta+4\sqrt e+4U/(U+t))\), where \(U\)
is the upper source deviation and \(\delta_{\rm row}\) is the
missing ambient-row proportion. The floor cuts this integral off
at \(HH\). The last fraction integrates to
\(4U\log(1+HH/U)\), taken as zero when \(U=0\). For \(t\ge HH\)
the original floor makes the strict lower tail empty. This is how
the target minimum, including on small quotients, enters the
original mean-slope estimate.
\src{src/Formalism/SharpRestriction/GlobalSchurBounds/TopDensityNGeSixteen.lean}{519}{same full and descended word maps}
\src{src/Formalism/SharpRestriction/GlobalSchurBounds/TopDensityNGeSixteen.lean}{639}{the matched exponent in the frame bound}
\src{src/Formalism/HarderNarasimhan/HarderNarasimhanProofs/TheoremFSharp.lean}{308}{integration of the target tail}.


The ramified comparison gives two conclusions after the same
compression calculation. Its first conclusion includes
\(\log|\mathrm{Word}|\) in the mean-slope loss. Its second conclusion
omits that term by applying the sharp lattice theorem.
The adapter \code{joint\_a\_ramified\_holds} selects the first
conclusion. Selecting the second gives the sharper bound at the
same target floor and on the same matrices.
\src{src/Formalism/SizeWindowBrackets/ChainSlopeWindows/JointDModelAppliesUnitLattice.lean}{236}{the two conclusions and the choice of the first}.

The explicit exponents distinguish the target height from the
frame errors. In the original component comparison put
\(M_0=\max(C_t,5)+3\). The source has
\(E_d=9M_0+1\), the gauge majorant has exponent
\(E_c=9M_0+2\), and the final exponent is at least
\(1+\max(E_b,E_c,C_{F0},E_d)\). In particular this choice is at
least \(75\). The raw matrix estimates give square-root bounds
\(\sqrt{d_e},\sqrt{e_0}\le\exp(-x^{\varepsilon/8})\)
before replacing them by the polynomial majorants. The factor
\(x^{C_H}\) in the height bound therefore pairs with
\(x^{-C_H}\) in the assembled frame transfer. The raw bounds control the square roots directly.
\src{src/Formalism/SharpRestriction/Divisibility/SlopeExcess/HeightBoundOfComponentBounds.lean}{73}{the component exponents}
\src{src/Formalism/SharpRestriction/Divisibility/HookProducts/GaugeFormBoundOfUniformSize.lean}{39}{the gauge exponent}
\src{src/Formalism/SharpRestriction/GlobalSchurBounds/TopDensityNGeSixteen.lean}{740}{the raw estimates before the polynomial weakening}.



The code-size condition uses this exponent too. With full word
rank \(r_f\), the original word-size bounds yield
\(D_*/\sigma\le r_f\) from
\(C_0\Gamma^2x^{C_H}\le|\mathrm{Word}|\).
Thus increasing \(C_H\) to accommodate a different floor also
requires checking the code size at that exponent.
\src{src/Formalism/SharpRestriction/GlobalSchurBounds/Chosen.lean}{56}{code size at the frame exponent}
\src{src/Formalism/SharpRestriction/GlobalSchurBounds/TopDensityNGeSixteen.lean}{549}{full rank from the code size and word ranks}.

The elementary part of the source-side budget can be isolated.
The original \code{ga\_pow} and \code{fo\_pow} apply on the
actual filled host, before the row separation or the
complement-two bound on \(-E_i\). They use the local cutoff,
the prescribed exponents, the Mertens bound, and the family size
bounds. They give \(G\le x^{4M_0+1}b\) and
\(F^0\le x^{9M_0+1}b\). Since \(x\ge2\) and \(M_0\ge8\), addition
gives \(G+F^0\le x^{9M_0+2}b\). It remains to bound the
signed \( -E_i\) contribution and the other source terms when
applying this estimate in a replacement for the full height proof.
\src{src/Formalism/SharpRestriction/Divisibility/HookProducts/CardLiftsEq.lean}{272}{the gain bound on the filled host}
\src{src/Formalism/SharpRestriction/Divisibility/HookProducts/LwPNonneg.lean}{250}{the forced-exponent bound}.

\paragraph{What the height comparison would need from other tests.}
The displayed implications also give a sufficient replacement
criterion independent of the exponent-grid proof. Suppose there is
an injective \(K\)-linear map \(\Psi\) on the primal row space, a
finite fractional ideal ledger with
\(\Psi(R_{\mathrm b})\subseteq (\prod_w\mathfrak p_w^{-a_w})
 M_{\mathrm{aux}}\), and an operator-height estimate in the
\(\kappa\)-normalized source metric. By
\eqref{geom:auxiliary-row-height} it suffices to bound the sum of the
auxiliary maximum slope, the operator height and the finite ideal
cost by \(x^A b\) for some fixed \(A\), with the original source and
passage bounds. For example a single positive rational integer
\(q_{\mathrm{den}}\) clearing the entire image into
\(M_{\mathrm{aux}}\) permits the finite cost \(2\log q_{\mathrm{den}}\):
the valuation of \(q_{\mathrm{den}}\mathcal O_K\) gives
\(\sum_w\lambda_wv_w(q_{\mathrm{den}})=\log q_{\mathrm{den}}\).
This deduction uses elementary ideal norms and the slope
inequality. The original source certifies its actual auxiliary
map through the fitted-grid and row-tower inputs above.
Producing a bounded such map on a new host from its elementary
host data alone is an additional geometric assertion.

\paragraph{A restricted test after the change of level.}
The stated \code{Passage} bounds the maximum and mean slopes. To
control an individual sharp-level subspace one can instead use
the common ambient lattice in the proof of that passage. Let
\(\Lambda\) denote this \(G\)-stable lattice and let \(e_0,e_1\)
be the base and sharp idempotents. Let \(h_0,h_1\) be their
subgroup orders. The embeddings
\(\iota_b,\iota_s\) into the ambient representation identify
the two row fibres with \(\operatorname{im}e_0,\operatorname{im}e_1\).
Their squared norms use the ambient metric divided by \(h_0,h_1\),
respectively. The same \(\kappa\) normalization on the two row
fibres preserves this ratio. The original sandwich gives
\[
 \iota_s(R_0)\subseteq\Lambda,\qquad
 h_0(\Lambda\cap\operatorname{im}e_0)\subseteq\iota_b(R_{\mathrm b}).
\]
Write \(\rho\) for the ambient group action and put
\(p_g=\iota_b^{-1}e_0\rho(g)\iota_s\). The maps \(p_g\) jointly
separate the entire sharp row fibre. Indeed the common kernel of
\(e_0\rho(g)\) in the ambient representation is \(G\)-stable. A
nonzero such kernel would contain an irreducible constituent.
The good projector is the identity and the base character
occurs once in each good irreducible. Thus \(e_0\) acts nontrivially
on that constituent, contradicting the \(g=1\) kernel condition.
These are the hypotheses used in the original level transfer.
\src{src/Formalism/LocalToGlobal/IndefiniteForms/IndefSide.lean}{339}{scope of the stated passage}
\src{src/Formalism/LocalToGlobal/GlobalTheory/DivPassage.lean}{70}{ambient metric and lattice sandwich}
\src{src/Formalism/LocalToGlobal/IndefiniteReal/LevelAmbTransfer.lean}{13}{the good projector and the base multiplicity}.

For a nonzero rational sharp row subspace \(A\) of rank \(a\)
choose a set \(C\subseteq G\) of at most \(a\) elements for which the
\(p_g|_A\) jointly separate \(A\). To construct it, start with the
whole \(A\) as kernel and add a \(g\) nonzero on the current
kernel until its dimension is zero. Set
\(A_b=\sum_{g\in C}p_g(A)\), so \(\dim A_b\le a^2\).
Suppose one has a fitted test \(\Psi_b\) injective on \(A_b\), with
finite containing ideal \(J\) for \(\Psi_b(R_{\mathrm b}\cap A_b)\)
in \(M_{\mathrm{aux}}\), and operator height at most \(h_\infty\)
on \(A_b\) in the normalized base metric. Define
\[
          J_A(y)=(\Psi_b p_g(y))_{g\in C}.
\]
This is injective on \(A\). Since \(h_0e_0\) is an integral group sum and \(\Lambda\) is \(G\)-stable,
\(h_0p_g\) maps \(R_0\cap A\) to the inverse image of
\(\Lambda\cap\operatorname{im}e_0\). The sandwich supplies the
second factor \(h_0\), so
\[
 h_0^2J_A(R_0\cap A)\subseteq
                  J\,M_{\mathrm{aux}}^{\oplus C}.
\]
At infinity \(e_0\) is a contraction, as it is an average of
unitary operators with unit-modulus coefficients. Each \(p_g\)
therefore has norm at most \(\sqrt{h_1/h_0}\). Taking the orthogonal
sum gives a squared operator cost of at most \(|C|h_1/h_0\) times
the squared cost of the base test. The extra height and ideal
allowance together are consequently
\[
       \log|C|+\log h_1+3\log h_0.
\]
For the top exterior line of \(A\) the allowance is at most \(a\)
times this amount. The auxiliary maximum slope of the orthogonal
sum is the same as that of \(M_{\mathrm{aux}}\).
This argument is a direct restriction of the underlying arithmetic
representation. It provides a test route for the designated
subspace without using any identification of base and sharp HN
cutoffs. To improve on the whole-space height bound, the actual
restricted operator or finite ideal cost still has to improve.
\src{src/Formalism/LocalToGlobal/IndefiniteReal/TransferStatement.lean}{30}{the integral multiple of the base projector}
The maximum-slope assertion for an orthogonal sum also follows
by filtering any sublattice by its successive coordinate kernels.
Each successive quotient maps integrally and contractively into
one coordinate lattice. Its slope is at most that coordinate's
maximum slope by the determinant inequality. Adding their
degrees gives the assertion.

\paragraph{The test lattice and the restricted infinite-place bound.}
The target of the fitted test has a separate elementary estimate.
The original \code{IsTestLat} uses the dual of the orthogonal
sum of the canonical class-representative coefficient lattices,
with the inverse-transpose Gram form. For this lattice
\code{leafMM\_of\_test} gives
\[
 \mu_{\max}(M_{\mathrm{aux}})
     \le 4f\,n_pR_p\log(n_pR_p),\qquad
 R_p=(m+1)(g_a+1)+L_c+n_p .
\]
The argument uses the self-dual lattices at the class
representatives and the coefficient-weight spread. It uses no
nonvanishing property of the fitted functional. On the actual
sizes \(f,k,m,t\) are bounded by a fixed power of \(x\);
\(n_p\) and the auxiliary rank \(N\) are polynomial expressions
in these ranks. Also \(L_c\) is the recorded cloud length, at
most \(\mathrm{Sc}\le xb/5\). The largest term in \(g_a\) involving
\(L_c\) is a polynomial in the ranks times \(\sqrt{L_c}\).
It follows that the displayed bound is
\(O(x^B\sqrt b)=o(b)\) for a fixed \(B\), because
\(b=\exp(x^\varepsilon)\). This last estimate is a consequence
of the numerical formulas and applies equally on the canonical
trivial-character host.
\src{src/Development/AbelianVarietyTheory/Definitions/Divisors/GridTestLatticeLargestSlopeBound.lean}{116}{the actual auxiliary lattice and metric}
\src{src/Development/AbelianVarietyTheory/Proofs/DivisorProofs/MultiplicityAtFittedPoints.lean}{94}{the bound for the canonical test lattice}
\src{src/Development/AbelianVarietyTheory/Statements/PlaceShiftStatements/SetAtOfSizeAt.lean}{60}{the cloud-length bound}
\src{src/Formalism/LocalRingCardinality/ColumnWeightSpreadOffTop.lean}{73}{the scale bound on the spread}.

The other infinite-place input has two steps. The fitted
\code{InfAt} estimate bounds the height of the ambient evaluation
map in the test metric. The quotient coisometry preserves its
operator norm on passage to the base quotient. Finally
\code{ip\_first} uses the lower type factors \(\lambda_j\) to
pass from that quotient metric to the normalized row metric.
In the squared-norm convention, if \(g(v)\) is a common lower
bound for \(\sum_j\log\lambda_j\) on the active types, then the
quotient operator bound gives the normalized row bound with
correction \(-\operatorname{mean}_v g(v)\).

This last comparison can be restricted explicitly to \(A_b\).
At each infinite place let \(Q_A=\iota_{\mu_0}^{-1}A_b\) in
base quotient coordinates and let \(g_A(v)\) be the minimum
of \(\sum_j\log\lambda_j\) over only the orthogonal types whose
projection of \(Q_A\) is nonzero. For \(q\in Q_A\), the norm
comparisons on each type and the orthogonal decomposition give
\[
 \|\iota_{\mu_0}q\|_{\kappa\text{-normalized},v}^{\,2}
       \ge e^{g_A(v)}\|q\|_{Q,v}^{\,2}.
\]
Applying the same calculation as for the full operator bounds
the restricted height by the height in the quotient metric
on \(Q_A\), minus \(\operatorname{mean}_v g_A(v)\).
This permits a better \(\lambda_j\) penalty if the selected
subspace avoids the worst types. The construction of the fitted
affine coordinates and the bound on their actual evaluation
operator remain necessary to make that improvement quantitative.
The trivial Dirichlet character in the final prime comparison
does not fix these automorphic types or their theta-test
normalizations.
\src{src/Development/AbelianVarietyTheory/Statements/PlaceShiftStatements/InfClassAt.lean}{176}{height of the fitted ambient evaluation operator}
\src{src/Development/AutomorphicArithmetic/ThetaCorrespondence/ThetaPairLifts/ThirdPlaceSeries/IpaTie.lean}{22}{coisometry and the fitted height bound}
\src{src/Formalism/LocalToGlobal/IndefiniteForms/IndefiniteTree/FunctionalLowerBoundOfConditionsNumeric.lean}{826}{norm comparison using the lower type factors}.

